\section{Arbitrary allocation profiles}
\label{sec:allocation89}
The largest permitted block does not determine how that resource is
actually used.  We retain the entire list of block sizes.  Throughout
this section the groups are complete probe--reference systems, not
merely probe marginals.  We first specify independent acquisition,
then allow the reset-constrained feedback of Definition~\ref{def:reset88}
while keeping the complete reservation profile fixed.

\begin{definition}[A prescribed allocation]\label{def:allocation89}
Let $\boldsymbol n=(n_1,\ldots,n_\ell)$ be a nonempty finite list of
positive integers.  Write
\[
 N_{\boldsymbol n}=\sum_{r=1}^{\ell}n_r,
 \qquad V_{\boldsymbol n}=\sum_{r=1}^{\ell}n_r^2.
\]
An allocation tester has an input
$\bigotimes_r\sigma_r$, where $\sigma_r$ is an arbitrary state on
$(\C^d)^{\otimes n_r}\otimes R_r$.  Reference dimensions and final
joint processing are unrestricted.  Public mixtures, with the
mixture label retained, are allowed for this fixed allocation.
A group may use fewer than its assigned slots by padding with fixed
inputs and discarding their outputs.  Denote the supremum of the
unhalved final trace separation by $D^{\boldsymbol n}(E,F)$.
For the same list, a prescribed-profile reset tester reserves $n_r$
calls in its $r$th fresh block, independent of the classical history.
The preparation and all internal adaptive controls may depend on that
history, but the complete fresh system is conditionally independent
of all old receiver memory.  This is one common conditional rule on
every formal history, not an equality of the two conditional receiver
states.  Early termination can be padded by unused reserved blocks.
Write $D^{\mathrm{reset},\boldsymbol n}$ for this larger supremum.
Both $N_{\boldsymbol n}$ and $V_{\boldsymbol n}$ are hard budgets,
not expectations.  Then
\[
 D^{\boldsymbol n}\le D^{\mathrm{reset},\boldsymbol n}\le D_{N_{\boldsymbol n}}.
\]
\end{definition}

The first moment $N_{\boldsymbol n}$ governs independent statistical
accumulation; the second moment $V_{\boldsymbol n}$ governs coherent
accumulation within these groups.  The second-moment principle has
established metrological antecedents \cite{Hyllus78,Toth78}.  The point
here is a finite trace-distance statement uniform over a whole
measurement neighborhood, including its unknown quadratic remainder.
We first record the finite classical estimate that permits unequal
blocks without grouping them into dyadic size classes.

\begin{lemma}[A common readout for unequal signals]\label{lem:unequalsignals89}
Let $0<\alpha\le1$, $0\le x_r\le1/2$, and
$\alpha x_r\le u_r\le x_r$.  Let $P_0$ be the law of independent
fair signs $X_r\in\{-1,1\}$, and let $P_u$ be the independent law
with $P_u\{X_r=1\}=1/2+u_r$.  There is a binary randomized readout
depending only on $(x_r)$, and not on $(u_r)$, for which the unhalved
separation is at least
\begin{equation}\label{eq:unequalsignals89}
 \frac{\alpha}{64}\min\left\{1,
                         \left(\sum_r x_r^2\right)^{1/2}\right\}.
\end{equation}
Consequently the full product laws have at least this separation.
\end{lemma}
\begin{proof}
Put $S=\sum_r x_r^2$.  For $S=0$ the assertion is immediate.
Otherwise set
\[
 A=\max\{\sqrt S,S\},\qquad
 t_r=\frac{x_r}{16A},\qquad
 T(X)=\sin\left(\sum_r t_rX_r\right).
\]
Then $0\le t_r\le1/16$ and $\sum_r t_r^2\le1/256$.
Under $P_0$ the expectation of $T$ is zero.  Independence under $P_u$
gives
\[
 \mathbb E_u\exp\left(i\sum_r t_rX_r\right)
 =\prod_r(\cos t_r+2iu_r\sin t_r)=R e^{i\phi},
 \quad \phi=\sum_r\arctan(2u_r\tan t_r).
\]
Each factor has positive real part.  Also
\[
 R\ge\prod_r\cos t_r
 \ge1-\tfrac12\sum_r t_r^2\ge\tfrac12.
\]
For $0\le t\le1/16$, $t\le\tan t\le2t$.  The arguments of
$\arctan$ above lie in $[0,1]$, where $v/2\le\arctan v\le v$.
It follows that
\[
 \frac{\alpha S}{16A}\le\sum_r u_rt_r\le\phi
 \le4\sum_r u_rt_r\le\frac{S}{4A}\le\frac14.
\]
Since $\sin\phi\ge\phi/2$ on this interval,
$\mathbb E_uT=R\sin\phi\ge\alpha S/(64A)$.
Accept with probability $(1+T(X))/2$.  The resulting binary laws
have unhalved separation $|\mathbb E_uT-\mathbb E_0T|$.
Finally $S/A=\min\{1,\sqrt S\}$.  This proves the assertion with
one readout for every permitted vector of biases.
\end{proof}

\begin{theorem}[Allocation-profile law]\label{thm:allocation89}
Fix an ordered measurement $E$, a nonzero Hermitian tuple $H$ with
$\sum_jH_j=0$ and $Q_jH_jQ_j\succeq0$, and a number $\Lambda\ge0$.
Here $P_j=\operatorname{supp}E_j$, $Q_j=I-P_j$, and
$\|R\|_\Sigma=\sum_j\|R_j\|_{\mathrm{op}}$.
There are $c,C,s_0>0$, depending on $E,H,\Lambda$ alone, such that
for every allocation $\boldsymbol n$, every $0<s\le s_0$, and every
legal measurement $F$ satisfying
\begin{equation}\label{eq:allocationtube89}
 \|F-E-sH\|_\Sigma\le\Lambda s^2,
\end{equation}
one has
\begin{equation}\label{eq:allocationlaw89}
 c\,\omega_{\boldsymbol n}(s)
 \le D^{\boldsymbol n}(E,F)
 \le D^{\mathrm{reset},\boldsymbol n}(E,F)\le C\,\omega_{\boldsymbol n}(s),
\end{equation}
where
\[
 \omega_{\boldsymbol n}(s)=
 \begin{cases}
 \min\{1,\sqrt{N_{\boldsymbol n}}s\},
                    &H\in\operatorname{Ran}\mathcal C_E,\\
 \min\{1,\sqrt{V_{\boldsymbol n}}s\},
                    &Q_jH_jQ_j=0\ (1\le j\le k),\quad
                           \Gamma_E(H)\notin\mathcal W_E,\\
 \min\{1,N_{\boldsymbol n}s\},&\text{some }Q_jH_jQ_j\ne0.
 \end{cases}
\]
These constants are uniform in the number and sizes of groups, and
in the entire permitted remainder.  In every row a single lower
tester can be chosen from $E,H,\Lambda,s,\boldsymbol n$, independently
of that remainder.  The coherent construction uses ideal
known-direction controls.  It is neither a common learner nor a
statement about separable probe marginals with a shared reference.
Nonemptiness and a sufficient allowance are as in
Proposition~\ref{prop:nonempty87}; its legality interval $s_*$ does
not replace the discrimination interval $s_0$.
\end{theorem}
\begin{proof}
Write $N=N_{\boldsymbol n}$ and $V=V_{\boldsymbol n}$.
Product-reference inputs are admissible in every allocation, and every prescribed-profile reset tester is an adaptive tester.  The regular and opening
rows therefore follow from Theorem~\ref{thm:tubedichotomy86}.
In the regular row its finite upper is
\[
 2\beta_h\sqrt N s+(\Lambda+K_h)Ns^2,
\]
capped at two, for the horizontal factor solution.  This is of the
asserted order since $x^2\le x$ for $x=\sqrt N s\le1$.
The opening lower uses the fixed impossible-output event, not a
coherent code.

First consider the independent coherent upper.  Use the canonical, generally nonhorizontal,
factor and the constants $\beta_0,r_0$ of
Lemma~\ref{lem:blockbures87}, writing $r_0$ for its remainder constant $r$.
With root fidelity
$f(\rho,\sigma)=\|\sqrt\rho\sqrt\sigma\|_1$ and Bures distance
$\beta^2=2(1-f)$, the normalization is
\[
 \beta(\rho,\sigma)^2\le\|\rho-\sigma\|_1
                         \le2\beta(\rho,\sigma).
\]
The actual channel remainder in a group of $n_r$ calls is
$n_r r_0s^2$; its Bures contribution is $s\sqrt{r_0n_r}$.
The block estimate consequently gives the finite upper
\begin{align}
 D^{\boldsymbol n}(E,F)
 &\le \min\left\{2,\,
  2s\left[\sum_r
      (\tfrac32\beta_0n_r+\sqrt{r_0n_r})^2\right]^{1/2}\right\}
                                                   \label{eq:profilefinite89}\\
 &\le\min\{2,\,2s(\tfrac32\beta_0\sqrt V+\sqrt{r_0N})\}.
                                                   \label{eq:profileminkowski89}
\end{align}
The last line is the Euclidean triangle inequality.  Since $N\le V$,
it implies the required upper.  Fidelity multiplies only between
independent group outputs.  A retained public mixture gives a direct
sum with common weights, whose trace separation is the weighted sum
of the conditional separations; the same bound follows.  Final
joint processing is unrestricted and contracts trace norm.

The same finite upper holds for prescribed-profile reset testers,
for a different reason.  Lemma~\ref{lem:adaptiveblock88} gives the
normalized fresh output deficit
\[
 a_r=\tfrac12s^2(\alpha n_r+\sqrt{r_0n_r})^2,
 \qquad \alpha=\tfrac32\beta_0,
\]
uniformly at every history in round $r$.  These deficits have the
same deterministic sum along every full reservation path.  Apply
Lemma~\ref{lem:resetfidelity88} to positive subnormalized receiver
operators.  If $A=\sum_ra_r\le1$, then
$\|\rho_E-\rho_F\|_1\le2\sqrt{2A}$; if $A>1$, use the trace cap two.
This proves \eqref{eq:profilefinite89} and
\eqref{eq:profileminkowski89} with $D^{\mathrm{reset},\boldsymbol n}$
in place of $D^{\boldsymbol n}$.  It does not multiply unconditional
output fidelities after feedback.  Countable histories are interpreted
as trace-class direct sums; all fidelity summands are nonnegative.

For the lower let $\Delta,\kappa>0$ be as in
Lemma~\ref{lem:parallelcode87}, decreasing $s_0$ if necessary so that
\[
 s_0\le\frac1{4\Delta},\qquad
 s_0\le\frac{\Delta}{\pi\kappa}.
\]
In group $r$ use
\[
 m_r=\min\{n_r,\lfloor(2\Delta s)^{-1}\rfloor\},\qquad
 x_r=m_r\Delta s,
\]
logical GHZ inputs, with the remaining slots padded and discarded.
All encodings precede acquisition and all recoveries follow it.
The reference dimension is at most $2d$ per used call; the reference
of a whole group may grow with its size.
The fixed binary event of that lemma has base probability $1/2$ and
alternative probability $1/2+u_r$, with
\[
 \left|u_r-\tfrac12\sin x_r\right|
       \le\tfrac12m_r\kappa s^2.
\]
This is the full finite tensor-channel error, paid before statistical
amplification.  Since $0<x_r\le1/2$, the displayed choice of $s_0$
implies
\[
 \frac{x_r}{2\pi}\le u_r
 \le\frac{x_r}{2}+\frac{\kappa s}{2\Delta}x_r\le x_r.
\]
Independent group outputs need not have equal biases.
Lemma~\ref{lem:unequalsignals89}, with $\alpha=1/(2\pi)$, gives
\begin{equation}\label{eq:profilelowerfinite89}
 D^{\boldsymbol n}(E,F)\ge
 \frac1{128\pi}\min\left\{1,
                      \Delta s\left(\sum_r m_r^2\right)^{1/2}\right\}.
\end{equation}
Both the individual events and the final randomized classical
readout are independent of $F-E-sH$.

If every $m_r=n_r$, the last minimum is
$\min\{1,\Delta s\sqrt V\}$, at least
$\min\{1,\Delta\}\min\{1,s\sqrt V\}$.
If some $m_r<n_r$, then $(2\Delta s)^{-1}\ge2$ and the floor bound
gives $x_r\ge1/4$.  The last minimum is consequently at least $1/4$.
The target $\min\{1,s\sqrt V\}$ is at most one, so this also controls
it by an absolute positive factor.  In particular the coherent lower
constant may be taken as $\min\{\Delta,1/4\}/(128\pi)$.
This proves the uniform lower for both classes, including profiles with very unequal
or very large groups.  No output is used to prepare a later input.
\end{proof}

\begin{corollary}[Width budgets and effective width]\label{cor:profilebudget89}
For $1\le b\le N$, put $q=\lfloor N/b\rfloor$ and $r=N-qb$.
Among allocations with $\sum n_j\le N$ and $n_j\le b$,
\begin{equation}\label{eq:profilemax89}
 \max V_{\boldsymbol n}=qb^2+r^2,
 \qquad \tfrac12Nb\le qb^2+r^2\le Nb.
\end{equation}
Thus the coherent width law can equivalently be written with
$\sqrt{qb^2+r^2}\,s$, with constants independent of $N,b$.
For a prescribed allocation the effective width is
$b_{\rm eff}=V_{\boldsymbol n}/N_{\boldsymbol n}$, not necessarily
its largest part.  Coarsening an allocation increases its coherent
scale and leaves the other two scales unchanged.
\end{corollary}
\begin{proof}
Adding unused slots increases the objective.  If two parts are below
$b$, transfer slots from the smaller to the larger until one is zero
or the other is $b$.  Their sum of squares does not decrease.
Iteration gives $q$ parts equal to $b$ and, if $r>0$, one part $r$.
The upper bound follows from $n_j^2\le bn_j$.
Since $r<b\le qb$, one has $qb\ge N/2$, which proves the lower.
Merging parts $u,v$ adds $2uv$ to their squared sum.  For independent allocation testers, the operational
inclusion also holds because product inputs within a merged group
remain admissible.  No analogous set inclusion is asserted for
arbitrary reset-feedback testers when a receiver-interaction boundary
is removed; their local orders are the preceding theorem's statement.
\end{proof}

For example, with $N=L^2$ calls and largest block $L$, the allocation
$(L,1,\ldots,1)$ has $V=2L^2-L$, whereas $L$ groups each of size $L$
have $V=L^3$.  Before saturation their coherent scales differ by a
factor of order $\sqrt L$, despite equal call counts and equal largest
blocks.  At the single-group endpoint, the preceding theorem recovers
matching parallel and adaptive \emph{local orders}; it makes no claim
of equality of their exact distances or global error exponents.

The comparison concerns the best tester for a prescribed reservation
profile.  It does not state that each feedback policy using that
profile is informative.  For example, a common final channel that
forgets every output has zero separation.  Likewise, an arbitrary
policy's worst-path cost is not itself a matching lower certificate:
a costly branch can be uninformative or rarely visited.
